% Conclusion Section

本文提出了一种基于事件驱动有限状态机的准被动双足步态控制方法，通过 V01-V24 迭代开发建立了从理论证明到硬件设计的完整技术路线。主要贡献包括：

\textbf{(1) 摩擦补偿的 FSM 控制律}：V13.7 证明前馈补偿 $\tau_{\text{comp}} = \tau_C \cdot \text{sign\_smooth}(\omega)$ 可完全消除 Dynamixel MX-64AR 电机的共振失败带（0.013-0.20 N·m），使准被动步态在真实电机约束下可行。

\textbf{(2) 传感器闭环状态估计}：V17 验证了低成本传感器（4096 CPR 编码器 + IMU）在 3× 噪声裕度下的鲁棒性，为硬件落地提供可行方案。

\textbf{(3) 侧向稳定性的理论边界}：V18-V20 否定结果揭示矢状面被动稳定与冠状面被动稳定的内在矛盾，为 2D/3D 混合控制策略提供物理依据。

\textbf{(4) 决策门方法论}：24 版本开发采用渐进式验证和决策门，避免盲目堆功能，确保技术积累可靠性。

\textbf{(5) 硬件设计完整方案}：V22 提供 2D 样机设计文档（机械结构、BOM、参数映射、系统辨识），物料成本 6500-8000 元，可直接用于制造。

\subsection{Future Work}

\textbf{短期（3 个月内）}：
\begin{itemize}
    \item \textbf{V25 实地行走}：完成 V22 硬件制造，进行真实地面行走测试
    \item \textbf{长时程仿真}：运行 100 s 验证 Floquet max$|\lambda| > 1.0$ 的长期影响
    \item \textbf{FSM 参数优化}：网格搜索使 max$|\lambda| < 0.95$，提升稳定裕度
\end{itemize}

\textbf{中期（6-12 个月）}：
\begin{itemize}
    \item \textbf{Sim-to-real 迁移}：基于 V23 系统辨识修正仿真模型，使用 V24 域随机化训练鲁棒控制器
    \item \textbf{多电机验证}：测试 EC-i 40（低摩擦）、AX-12A（低成本）的 $\tau_C$ 特性
    \item \textbf{能量优化}：引入自适应能量反馈，将 CoT 从 0.066 降至 $<0.05$（接近人类水平）
\end{itemize}

\textbf{长期（1-2 年）}：
\begin{itemize}
    \item \textbf{3D 硬件验证}：基于 V25 的硬件数据，用实测漂移率和执行能力重新设计 3D 控制策略
    \item \textbf{地形适应}：扩展到崎岖地形（V11 的硬件版本）
    \item \textbf{学习辅助}：用强化学习（RL）在线调整 FSM 参数，适应磨损和环境变化
\end{itemize}

\subsection{Broader Impact}

本研究的方法论（迭代开发 + 决策门 + 否定结果透明化）可推广到其他机器人系统：

\begin{itemize}
    \item \textbf{教育}：V01-V24 完整代码和文档可作为机器人课程的案例研究
    \item \textbf{低成本机器人}：V22 设计（<8000 元）降低准被动行走的研究门槛
    \item \textbf{开源社区}：公开否定结果（V18-V20）避免其他研究者重复失败路径
    \item \textbf{工业应用}：准被动步态的高能量效率适用于巡检、配送等长时间作业场景
\end{itemize}

\subsection{Closing Remarks}

准被动双足行走在能量效率和自然动力学利用上展现巨大潜力，但从仿真到硬件的迁移面临系统性挑战。本文通过 24 版本的渐进式验证，揭示了电机摩擦、传感器噪声和侧向稳定的物理极限，为硬件落地提供了可验证的设计方法论。未来工作将聚焦于 V25 实地行走验证和 3D 化路线的重新审视。我们相信，通过硬件实测数据驱动的算法迭代，准被动行走机器人可在实际应用中达到人类级别的能量效率和鲁棒性。
